\subsubsection{高斯消元 (Gaussian Elimination)}
高斯消元法，用于求解线性方程组。
\begin{lstlisting}
double a[110][110];
double b[110];
ll p[110];
const double eps=1e-7;
int main(){
	n=d;f(i,1,n){
		f(j,1,n+1)cin>>a[i][j];
	}ll row=1;
    f(i,1,n){
		ll pos=row;f(j,row,n)if(fabs(a[j][i])>fabs(a[pos][i]))pos=j;
        if(row>n||fabs(a[pos][i])<eps)continue;
		swap(a[row],a[pos]);
		double x=a[row][i];f(j,1,n+1)a[row][j]/=x;
		f(j,row+1,n){
			x=a[j][i];
			f(k,i,n+1)a[j][k]-=x*a[row][k];
		}row++;
	}
    // f(i,1,n){
    //     f(j,1,n+1)cout<<a[i][j]<<" ";cout<<endl;
    // }
    f(i,1,n){
        bool fl=1;
        f(j,1,n)if(fabs(a[i][j])>eps)fl=0;
        if(fl){if(fabs(a[i][n+1])>eps){puts("-1");return 0;}}//无解
    }
    f(i,1,n){
        bool fl=1;
        f(j,1,n)if(fabs(a[i][j])>eps)fl=0;
        if(fl){if(fabs(a[i][n+1])<=eps){puts("0");return 0;}}//无限解
    }
    for(int i=n;i>=1;i--){
		b[i]=a[i][n+1];
		f(j,1,i-1)a[j][n+1]-=a[i][n+1]*a[j][i];
	}
    f(i,1,n)printf("x%lld=%.2lf\n",(long long)i,b[i]);
	return 0;
}
\end{lstlisting}

\paragraph{多测与注意} 每组重读增广矩阵，row=1；b 会覆盖。消元破坏输入，重复使用先复制。eps 按题目精度调整。
